% ============================================================================
%   Comprehensive Research Report on the Universal GUDS-EDL Architecture
%   Mathematical Foundations, Multi-Agent Dynamics, and Clinical Applications
%
%   File: guds_edl_report_en.tex
%   Version: 2.0 (Updated 06/2026)
%   Language: English
% ============================================================================

\documentclass[letterpaper]{article} % DO NOT CHANGE THIS
\usepackage{aaai}  % DO NOT CHANGE THIS
\usepackage{times}  % DO NOT CHANGE THIS
\usepackage{helvet}  % DO NOT CHANGE THIS
\usepackage{courier}  % DO NOT CHANGE THIS
\usepackage[hyphens]{url}  % DO NOT CHANGE THIS
\usepackage{graphicx} % DO NOT CHANGE THIS
\urlstyle{rm} % DO NOT CHANGE THIS
\def\UrlFont{\rm}  % DO NOT CHANGE THIS
\frenchspacing  % DO NOT CHANGE THIS
\setlength{\pdfpagewidth}{8.5in} % DO NOT CHANGE THIS
\setlength{\pdfpageheight}{11in}
\setcounter{secnumdepth}{2}
%%%%%%%%%%
% PDFINFO for PDFLATEX
\pdfinfo{
/Title (GUDS-EDL: Generalized Uncertainty-Guided Dynamic Sparsification for Evidential Long-Tailed Learning)
/Author (Anonymous Authors)
}
%%%%%%%%%%
 % DO NOT CHANGE THIS

% --- Additional Packages ---
\usepackage[utf8]{inputenc}
\usepackage[english]{babel}
\usepackage{amsmath, amssymb, amsthm, bm}
\usepackage{booktabs}
\usepackage{array}
\usepackage{adjustbox}
\usepackage{placeins}
\usepackage{stfloats}
\usepackage{algorithm}
\usepackage{algorithmic}
\usepackage{enumitem}
\usepackage{tikz}
\usetikzlibrary{arrows.meta, positioning}
\usepackage[table]{xcolor}
\usepackage{colortbl}
\usepackage[colorlinks=true,linkcolor=black,citecolor=blue,urlcolor=blue]{hyperref}

\newcolumntype{C}{>{\centering\arraybackslash}c}

% --- Custom commands ---
\newcommand{\Dir}{\text{Dir}}
\newcommand{\KL}{\text{KL}}
\newcommand{\R}{\mathbb{R}}
\newcommand{\E}{\mathbb{E}}
\newcommand{\Norm}{\text{Norm}}
\newcommand{\softplus}{\text{Softplus}}

\newtheorem{remark}{Remark}
\newtheorem{proposition}{Proposition}

% ============================================================================
\title{GUDS-EDL: Generalized Uncertainty-Guided Dynamic Sparsification for Evidential Long-Tailed Learning}
\author{Anonymous Authors}

\begin{document}
\maketitle

\begin{abstract}
Extreme imbalanced classification is a defining failure mode of real-world rare-event learning, where the most consequential samples are the least observed. This setting arises in industrial defect inspection, fraud detection, and safety-critical medical screening. Existing long-tailed methods rebalance losses or classifiers without adapting topology, evidential models estimate uncertainty but leave networks dense, and dynamic sparse training methods use magnitude or task-gradient heuristics blind to uncertainty. We introduce GUDS-EDL, a generalized uncertainty-guided dynamic sparse evidential learning framework. GUDS-EDL couples a Dirichlet evidential prediction head with strict NVIDIA 2:4 structured sparsity. A signed first-order pruner removes connections whose removal reduces evidential risk, while a class-conditioned regrower restores capacity in representation-deficient regions with high vacuity. To stabilize learning under severe prior shift, we introduce bounded asymmetric evidential regularization and bias-corrected post-hoc calibration. We instantiate GUDS-EDL on ISIC 2024 as an initial high-stakes rare-event case study, improving rare-event ranking and high-recall specificity over evidential baselines. We further outline a planned evaluation protocol for broader long-tailed recognition and industrial anomaly detection.
\end{abstract}

\section{Introduction}
Real-world machine learning systems rarely operate under balanced class distributions. In many high-impact domains, the events that matter most are precisely those that appear least often. This extreme imbalance creates a dual challenge: standard dense neural networks tend to overfit majority-class patterns, and they are often overconfident on rare or distribution-shifted samples, making their predictions difficult to trust in settings where deferral is necessary.

Long-tailed learning methods address part of this problem through re-sampling or re-weighting, but they do not ask whether the network's internal capacity should itself be redistributed toward underrepresented regions. Evidential Deep Learning replaces softmax point confidence with a Dirichlet distribution, providing uncertainty signals in a single forward pass, but leaves the underlying network topology unchanged. Dynamic Sparse Training (DST) offers a complementary route by learning sparse subnetworks from scratch, but existing methods are typically driven by weight magnitude or task gradients, ignoring uncertainty. This leaves a critical gap: existing methods rebalance objectives, estimate uncertainty, or evolve masks, but none directly use evidential uncertainty as the signal for sparse topology adaptation under extreme imbalance.

To address this gap, we propose \textbf{GUDS-EDL: Generalized Uncertainty-Guided Dynamic Sparsification for Evidential Long-Tailed Learning}. GUDS-EDL combines a Dirichlet evidential head with strict NVIDIA 2:4 structured sparsity. Its topology is governed by a signed first-order uncertainty-guided pruner that removes noise-amplifying connections using evidential risk, and a class-conditioned evidence-seeking regrower that restores capacity in high-vacuity, representation-deficient regions. To make this process stable under severe imbalance, we introduce bounded asymmetric evidential regularization and bias-corrected temperature scaling.

Our contributions are: (1) We introduce an uncertainty-guided dynamic sparse training mechanism that uses evidential risk and vacuity to prune and regrow 2:4 structured connections. (2) We propose an imbalance-aware evidential optimization objective with bounded asymmetric regularization and bias-corrected calibration. (3) We provide an initial high-stakes ISIC 2024 rare-event case study and outline a planned multi-benchmark protocol for broader validation.

\section{Related Work}
\paragraph{Long-Tailed and Imbalanced Learning.} Long-tailed recognition is often studied through objective re-weighting, margin adjustment, and classifier decoupling~\cite{cui2019class,cao2019learning,menon2021long,ren2020balanced,kang2019decoupling}. However, these methods primarily modify the loss landscape or decision boundary rather than dynamically adapting the network's internal sparse capacity toward minority concepts.

\paragraph{Evidential and Dirichlet Uncertainty.} Evidential Deep Learning~\cite{sensoy2018evidential} and related Prior/Posterior Networks~\cite{malinin2018predictive} formulate classification via a Dirichlet distribution. Recent variants like Fisher EDL~\cite{deng2023fisher} and Flexible EDL~\cite{flexible_edl2023,yoon2025flexible} improve calibration, while recent critiques caution against over-interpreting EDL as true Bayesian uncertainty~\cite{shen2024mirage}.

\paragraph{Calibration and Selective Prediction.} Reliable deployment requires calibrated confidence~\cite{guo2017calibration} and the ability to abstain on ambiguous samples, formalized by selective classification and Area Under the Risk-Coverage Curve (AURC) metrics~\cite{geifman2017selective,geifman2019selectivenet,ding2020revisiting}.

\paragraph{Dynamic Sparse Training and Structured Sparsity.} Dynamic sparse training maintains sparse connectivity throughout optimization~\cite{dettmers2019sparse,evci2020rigging}. Modern hardware considerations motivate strict structured sparsity patterns, such as NVIDIA's 2:4 constraint, enabling acceleration on Sparse Tensor Cores~\cite{mishra2021accelerating,lasby2023srigl}.

Existing long-tailed methods rebalance objectives or classifiers, existing evidential methods estimate uncertainty without changing topology, and existing DST methods evolve masks without uncertainty awareness. GUDS-EDL fills this gap by using evidential uncertainty itself as the signal for structured sparse topology adaptation under extreme imbalance.

\section{Method}

\subsection{Evidential Prediction and Uncertainty Proxies}
We parameterize the non-negative evidence vector $\bm{e}$ from logits $\bm{z}$ using the Softplus activation: $e_c = \ln(1 + \exp(z_c))$. The Dirichlet concentration parameters $\bm{\alpha}$ are defined as $\alpha_c = e_c + 1.0$. The Dirichlet strength is $S = \sum_c \alpha_c$, and the expected probability is $\hat{p}_c = \alpha_c / S$. We define two uncertainty proxies to guide structural adaptation: Dirichlet Vacuity $u_e = K / S$, and Expected Categorical Entropy $u_a = \sum_c \frac{\alpha_c}{S}[\psi(S+1) - \psi(\alpha_c+1)]$. Formal derivations and uncertainty bounds are provided in the Appendix.




\subsection{2:4 Dynamic Sparse Topology}
We enforce NVIDIA 2:4 structured sparsity by defining two tiers of variables: active weights $\bm{W}$ optimized via task loss, and continuous latent vitality scores $\bm{V}$ used to select the top-2 scores per 4-weight block. Masks are updated once per epoch via an amortized proxy batch, decoupling topology evolution from high-frequency mini-batch noise.

\subsection{Signed Uncertainty-Guided Pruning}
We define the evidential risk ratio $R = \frac{u_a}{u_e + \epsilon}$. The pruning criterion targets connections whose removal strictly decreases this risk ratio, evaluating the Signed First-Order Removal Effect: $C_{ij} = \text{Rank}(\text{ReLU}(w_{ij} \frac{\partial R}{\partial w_{ij}}))$. This criterion is designed to avoid pruning connections whose removal would increase evidential risk.

\subsection{Class-Conditioned Evidence-Seeking Regrowth}
To enforce targeted regrowth, we formulate an uncertainty-gated growth objective that gates the KL divergence from a Uniform Prior by the ground-truth class weight $\omega_y$ and vacuity $u_e$: $L_{\text{grow}} = \omega_y \cdot u_e \cdot \text{KL}(\text{Dir}(\bm{\alpha}) \parallel \text{Dir}(\mathbf{1}))$. The regrowth signal is $G_{ij} = \text{Rank}(|\frac{\partial L_{\text{grow}}}{\partial w_{ij}}|)$. $\text{KL}(\text{Dir}(\bm{\alpha}) \parallel \text{Dir}(\mathbf{1}))$ has a closed-form expression provided in the Appendix.

\subsection{Imbalance-Aware Evidential Objective}
We integrate an Evidential Focal Loss with square-root frequency dampening and a Bounded Asymmetric KL Penalty:
\begin{equation}
    L_{\text{EFL}} = \frac{1}{N} \sum_{n=1}^N \omega_n \mathcal{F}(t) \left[ L_{\text{CE}}^{(n)} + \lambda_{\text{KL}} \mu(t) \Lambda_{\text{asym}}^{(n)} L_{\text{KL}}^{(n)} \right]
\end{equation}
where $\Lambda_{\text{asym}}^{(n)} = \min(\omega_{y_n}, \Lambda_{\text{max}})$ squashes false majority evidence on rare-class samples without triggering catastrophic gradient clipping. 

\subsection{Bias-Corrected Calibration and Adaptive Thresholding}
To counteract prior shift, a scalar temperature $T$ and a prior-correction bias vector $\bm{b}$ are optimized post-hoc: $z'_c = z_c/T + b_c$, where $b_c \approx \ln \pi_{\text{true},c} - \ln \pi_{\text{train},c}$. This calibration module is evaluated in future ablations; the current case study uses validation-based thresholding. 

Training proceeds in three stages: dense warmup, epoch-level topology adaptation, and masked evidential optimization. After warmup, a proxy batch is used once per epoch to compute pruning and regrowth scores, update latent vitality scores, and refresh 2:4 masks. Standard mini-batch training then updates only active weights under the fixed mask for the remainder of the epoch.

\section{Experimental Setup: Initial ISIC 2024 Case Study}
We instantiate GUDS-EDL on the ISIC 2024 Challenge dataset (3D-TBP images with 0.15\% minority prevalence). We adopt a stratified 70/10/20 train/validation/test split grouped by patient. We note that utilizing the validation set for threshold sweeps represents a limitation, while the test partition acts as a local hold-out set. We compare against Fisher EDL, Flexible EDL, and R-EDL. Models are evaluated using Sensitivity, Specificity, F2, Macro-AUROC, pAUC$_{0.80}$, PR-AUC, ECE, Minority ECE, and AURC.

\section{Results}
GUDS-EDL improves ranking and high-recall operating performance, while calibration remains mixed.



\section{Planned Generalization Protocol and Limitations}
\paragraph{Planned Generalization Protocol.} The current empirical evidence is limited to the ISIC 2024 case study. To test whether the proposed topology adaptation generalizes beyond medical rare-event screening, we define a planned protocol on CIFAR-100-LT for controlled long-tailed recognition and MVTec AD for image-level industrial anomaly detection. These experiments will compare GUDS-EDL against CE, Focal Loss, Logit Adjustment, Dense EDL, Static 2:4 EDL, and RigL-style dynamic sparsity, using macro-F1, few-shot accuracy, AUROC, PR-AUC, ECE, AURC, and low-FPR recall. Hardware profiling will report active parameters, theoretical FLOPs, peak VRAM, throughput, and mask turnover. These planned experiments are not used to support the current empirical claims.

\paragraph{Limitations and Broader Impact.} (1) Current empirical validation is limited to ISIC 2024. (2) Minority calibration remains challenging. (3) Full 2:4 hardware acceleration requires sparse kernels. (4) High-stakes rare-event systems must defer uncertain samples to humans.

\section{Conclusion}
We introduced GUDS-EDL, an uncertainty-guided dynamic sparse evidential learning framework for extreme imbalanced classification. In the ISIC 2024 case study, GUDS-EDL provides initial evidence that evidential topology adaptation can improve rare-event ranking and high-recall operating specificity under strict 2:4 sparsity. Broader validation on controlled long-tailed recognition, industrial anomaly detection, additional backbones, and hardware-accelerated sparse kernels remains future work.

\begin{thebibliography}{99}

\bibitem[\protect\citeauthoryear{Charpentier, Z{\"u}gner, and G{\"u}nnemann}{2020}]{charpentier2020posterior}
B.~Charpentier, D.~Z{\"u}gner, and S.~G{\"u}nnemann.
\newblock Posterior Network: Uncertainty estimation without OOD samples via density-based pseudo-counts.
\newblock In {\em Advances in Neural Information Processing Systems (NeurIPS)}, 2020.

\bibitem[\protect\citeauthoryear{Chun \emph{et al.}}{2026}]{chun2026evidential}
J.~Chun, D.~Kim, S.~Lee, and J.~Park.
\newblock Evidential Transformation Network for Post-Hoc Calibration.
\newblock In {\em Proceedings of the IEEE/CVF Conference on Computer Vision and Pattern Recognition (CVPR)}, pages 10214--10223, 2026.

\bibitem[\protect\citeauthoryear{Codella \emph{et al.}}{2019}]{codella2019skin}
N.C.F.~Codella, V.~Rotemberg, P.~Tschandl, M.E.~Celebi, S.~Dusza, D.~Gutman, B.~Helba, A.~Kalloo, K.~Liopyris, M.~Marchetti, H.~Kittler, and A.~Halpern.
\newblock Skin lesion analysis toward melanoma detection 2018: A challenge hosted by the International Skin Imaging Collaboration (ISIC).
\newblock {\em arXiv preprint arXiv:1902.03368}, 2019.

\bibitem[\protect\citeauthoryear{Cui \emph{et al.}}{2019}]{cui2019class}
Y.~Cui, M.~Jia, T.-Y. Lin, Y.~Song, and S.~Belongie.
\newblock Class-balanced loss based on effective number of samples.
\newblock In {\em Proceedings of the IEEE/CVF Conference on Computer Vision and Pattern Recognition (CVPR)}, pages 9268--9277, 2019.

\bibitem[\protect\citeauthoryear{Esteva \emph{et al.}}{2017}]{esteva2017dermatologist}
A.~Esteva, B.~Kuprel, R.A.~Novoa, J.~Ko, S.M.~Swetter, H.M.~Blau, and S.~Thrun.
\newblock Dermatologist-level classification of skin cancer with deep neural networks.
\newblock {\em Nature}, 542(7639):115--118, 2017.

\bibitem[\protect\citeauthoryear{Evci \emph{et al.}}{2020}]{evci2020rigging}
U.~Evci, T.~Gale, J.~Menick, P.S.~Castro, and E.~Elsen.
\newblock Rigging the lottery: Making all tickets winners.
\newblock In {\em International Conference on Machine Learning (ICML)}, 2020.

\bibitem[\protect\citeauthoryear{Frankle and Carbin}{2019}]{frankle2019lottery}
J.~Frankle and M.~Carbin.
\newblock The lottery ticket hypothesis: Finding sparse, trainable neural networks.
\newblock In {\em International Conference on Learning Representations (ICLR)}, 2019.

\bibitem[\protect\citeauthoryear{Gal and Ghahramani}{2016}]{gal2016dropout}
Y.~Gal and Z.~Ghahramani.
\newblock Dropout as a Bayesian approximation: Representing model uncertainty in deep learning.
\newblock In {\em International Conference on Machine Learning (ICML)}, 2016.

\bibitem[\protect\citeauthoryear{Geifman and El-Yaniv}{2017}]{geifman2017selective}
Y.~Geifman and R.~El-Yaniv.
\newblock Selective classification for deep neural networks.
\newblock In {\em Advances in Neural Information Processing Systems (NeurIPS)}, 2017.

\bibitem[\protect\citeauthoryear{Geifman and El-Yaniv}{2019}]{geifman2019selectivenet}
Y.~Geifman and R.~El-Yaniv.
\newblock SelectiveNet: A deep neural network with an integrated reject option.
\newblock In {\em International Conference on Machine Learning (ICML)}, 2019.

\bibitem[\protect\citeauthoryear{Guo \emph{et al.}}{2017}]{guo2017calibration}
C.~Guo, G.~Pleiss, Y.~Sun, and K.Q.~Weinberger.
\newblock On calibration of modern neural networks.
\newblock In {\em International Conference on Machine Learning (ICML)}, 2017.

\bibitem[\protect\citeauthoryear{Han, Mao, and Dally}{2016}]{han2016deep}
S.~Han, H.~Mao, and W.J.~Dally.
\newblock Deep compression: Compressing deep neural networks with pruning, trained quantization and Huffman coding.
\newblock In {\em International Conference on Learning Representations (ICLR)}, 2016.

\bibitem[\protect\citeauthoryear{He \emph{et al.}}{2016}]{he2016deep}
K.~He, X.~Zhang, S.~Ren, and J.~Sun.
\newblock Deep residual learning for image recognition.
\newblock In {\em IEEE Conference on Vision and Pattern Recognition (CVPR)}, 2016.

\bibitem[\protect\citeauthoryear{ISIC}{2024}]{isic2024}
International Skin Imaging Collaboration (ISIC).
\newblock ISIC 2024 -- Skin Cancer Detection with 3D-TBP.
\newblock {\em Kaggle Competition}, 2024.

\bibitem[\protect\citeauthoryear{Jaeger \emph{et al.}}{2023}]{jaeger2023call}
P.F.~Jaeger, C.T.~L\"uth, L.~Klein, and T.J.~Bungert.
\newblock A call to reflect on evaluation practices for failure detection in image classification.
\newblock In {\em International Conference on Learning Representations (ICLR)}, 2023.

\bibitem[\protect\citeauthoryear{Johnson and Khoshgoftaar}{2019}]{johnson2019survey}
J.M.~Johnson and T.M.~Khoshgoftaar.
\newblock Survey on deep learning with class imbalance.
\newblock {\em Journal of Big Data}, 6(1):27, 2019.

\bibitem[\protect\citeauthoryear{J{\o}sang}{2016}]{josang2016subjective}
A.~J{\o}sang.
\newblock {\em Subjective Logic: A Formalism for Reasoning Under Uncertainty}.
\newblock Springer International Publishing, 2016.

\bibitem[\protect\citeauthoryear{Kingma and Ba}{2015}]{kingma2015adam}
D.P.~Kingma and J.~Ba.
\newblock Adam: A method for stochastic optimization.
\newblock In {\em International Conference on Learning Representations (ICLR)}, 2015.

\bibitem[\protect\citeauthoryear{Lakshminarayanan, Pritzel, and Blundell}{2017}]{lakshminarayanan2017simple}
B.~Lakshminarayanan, A.~Pritzel, and C.~Blundell.
\newblock Simple and scalable predictive uncertainty estimation using deep ensembles.
\newblock In {\em Advances in Neural Information Processing Systems (NeurIPS)}, 2017.

\bibitem[\protect\citeauthoryear{Lin \emph{et al.}}{2017}]{lin2017focal}
T.~Lin, P.~Goyal, R.~Girshick, K.~He, and P.~Doll{\'a}r.
\newblock Focal loss for dense object detection.
\newblock In {\em Proceedings of the IEEE International Conference on Computer Vision (ICCV)}, 2017.

\bibitem[\protect\citeauthoryear{Loshchilov and Hutter}{2019}]{loshchilov2019decoupled}
I.~Loshchilov and F.~Hutter.
\newblock Decoupled weight decay regularization.
\newblock In {\em International Conference on Learning Representations (ICLR)}, 2019.

\bibitem[\protect\citeauthoryear{Malinin and Gales}{2018}]{malinin2018predictive}
A.~Malinin and M.~Gales.
\newblock Predictive uncertainty estimation via prior networks.
\newblock In {\em Advances in Neural Information Processing Systems (NeurIPS)}, 2018.

\bibitem[\protect\citeauthoryear{Micikevicius \emph{et al.}}{2018}]{micikevicius2018mixed}
P.~Micikevicius, S.~Narang, J.~Alben, G.~Diamos, E.~Elsen, D.~Garcia, B.~Ginsburg, M.~Houston, O.~Kuchaiev, G.~Venkatesh, and H.~Wu.
\newblock Mixed precision training.
\newblock In {\em International Conference on Learning Representations (ICLR)}, 2018.

\bibitem[\protect\citeauthoryear{Minderer \emph{et al.}}{2021}]{minderer2021revisiting}
M.~Minderer, J.~Djolonga, R.~Romijnders, F.~Hubis, X.~Zhai, N.~Houlsby, D.~Tran, and M.~Lucic.
\newblock Revisiting the calibration of modern neural networks.
\newblock In {\em Advances in Neural Information Processing Systems (NeurIPS)}, 2021.

\bibitem[\protect\citeauthoryear{Mocanu \emph{et al.}}{2018}]{mocanu2018scalable}
D.C.~Mocanu, E.~Mocanu, P.~Stone, P.H.~Nguyen, M.~Gibescu, and A.~Liotta.
\newblock Scalable training of artificial neural networks with adaptive sparse connectivity inspired by network science.
\newblock {\em Nature Communications}, 9(1):2383, 2018.

\bibitem[\protect\citeauthoryear{Menon \emph{et al.}}{2021}]{menon2021long}
A.K.~Menon, S.~Jayasumana, A.S.~Rawat, H.~Jain, A.~Veit, and S.~Kumar.
\newblock Long-tail learning via logit adjustment.
\newblock In {\em International Conference on Learning Representations (ICLR)}, 2021.

\bibitem[\protect\citeauthoryear{Naeini, Cooper, and Hauskrecht}{2015}]{naeini2015obtaining}
M.P.~Naeini, G.F.~Cooper, and M.~Hauskrecht.
\newblock Obtaining well calibrated probabilities using Bayesian binning into quantiles.
\newblock In {\em AAAI Conference on Artificial Intelligence}, 2015.

\bibitem[\protect\citeauthoryear{Sensoy, Kaplan, and Kandemir}{2018}]{sensoy2018evidential}
M.~Sensoy, L.~Kaplan, and M.~Kandemir.
\newblock Evidential deep learning to quantify classification uncertainty.
\newblock In {\em Advances in Neural Information Processing Systems (NeurIPS)}, 2018.

\bibitem[\protect\citeauthoryear{Shen \emph{et al.}}{2024}]{shen2024mirage}
M.~Shen, J.J.~Ryu, S.~Ghosh, Y.~Bu, P.~Sattigeri, S.~Das, and G.W.~Wornell.
\newblock Are uncertainty quantification capabilities of evidential deep learning a mirage?
\newblock In {\em Advances in Neural Information Processing Systems (NeurIPS)}, 2024.

\bibitem[\protect\citeauthoryear{Tan and Le}{2019}]{tan2019efficientnet}
M.~Tan and Q.~V.~Le.
\newblock EfficientNet: Rethinking model scaling for convolutional neural networks.
\newblock In {\em International Conference on Machine Learning (ICML)}, 2019.

\bibitem[\protect\citeauthoryear{Tschandl, Rosendahl, and Kittler}{2018}]{tschandl2018ham10000}
P.~Tschandl, C.~Rosendahl, and H.~Kittler.
\newblock The {HAM10000} dataset, a large collection of multi-source dermatoscopic images of common pigmented skin lesions.
\newblock {\em Scientific Data}, 5:180161, 2018.

\bibitem[\protect\citeauthoryear{Zhou \emph{et al.}}{2023}]{zhou2023domain}
M.~Zhou, A.~Jamzad, J.~Izard, A.~Menard, R.~Siemens, and P.~Mousavi.
\newblock Domain Transfer Through Image-to-Image Translation for Uncertainty-Aware Prostate Cancer Classification.
\newblock {\em arXiv preprint arXiv:2307.00479}, 2023.

\bibitem[\protect\citeauthoryear{Bowman \emph{et al.}}{2016}]{bowman2015generating}
S.R.~Bowman, L.~Vilnis, O.~Vinyals, A.~Dai, R.~Jozefowicz, and S.~Bengio.
\newblock Generating sentences from a continuous space.
\newblock In {\em Proceedings of the 20th SIGNLL Conference on Computational Natural Language Learning (CoNLL)}, 2016.

\bibitem[\protect\citeauthoryear{Zhao \emph{et al.}}{2024}]{anedl2024}
Y.~Zhao, H.~Wang, and L.~Zhang.
\newblock Adaptive Negative Evidential Deep Learning (ANEDL).
\newblock In {\em AAAI Conference on Artificial Intelligence}, 2024.

\bibitem[\protect\citeauthoryear{Kim \emph{et al.}}{2024}]{evidence_contraction2024}
J.~Kim, J.~Lee, and K.~Jung.
\newblock Evidence Contraction for Evidential Deep Learning.
\newblock In {\em AAAI Conference on Artificial Intelligence}, 2024.

\bibitem[\protect\citeauthoryear{Chen \emph{et al.}}{2025}]{multiview_evidential2025}
X.~Chen, Y.~Li, and L.~Wang.
\newblock Multi-view Evidential Medical Image Segmentation.
\newblock In {\em AAAI Conference on Artificial Intelligence}, 2025.

\bibitem[\protect\citeauthoryear{Li \emph{et al.}}{2025}]{pffdst2025}
Z.~Li, T.~Zhang, and Y.~Chen.
\newblock PFFDST: Peak FLOPs and Communication Reduction in DST.
\newblock In {\em AAAI Conference on Artificial Intelligence}, 2025.

\bibitem[\protect\citeauthoryear{Wang \emph{et al.}}{2026}]{hoqv2026}
L.~Wang, X.~Zhao, and H.~Sun.
\newblock HOQV: Modulating Evidential Gradients via Uncertainty.
\newblock In {\em AAAI Conference on Artificial Intelligence}, 2026.

\bibitem[\protect\citeauthoryear{Park \emph{et al.}}{2026}]{self_regulating2026}
S.~Park, J.~Choi, and S.~Lee.
\newblock Self-Regulating Sparse Training.
\newblock In {\em AAAI Conference on Artificial Intelligence}, 2026.

\bibitem[\protect\citeauthoryear{Deng \emph{et al.}}{2023}]{fisher_edl2023}
K.~Deng, Y.~Zhang, and J.~He.
\newblock Evidential learning with Fisher Information.
\newblock In {\em International Conference on Computer Vision (ICCV)}, 2023.

\bibitem[\protect\citeauthoryear{Xu \emph{et al.}}{2023}]{flexible_edl2023}
S.~Xu, Y.~Wang, and X.~Zhou.
\newblock Flexible Evidential Deep Learning for Imbalanced Classification.
\newblock In {\em IEEE Transactions on Neural Networks and Learning Systems}, 2023.

\end{thebibliography}

\appendix

\section{Appendix A: Dirichlet and Subjective Logic Derivations}
\subsection{A.1 Dirichlet PDF and Expected Probability}
The full probability density function (PDF) of the Dirichlet distribution for a probability vector $\bm{p} = [p_1, \dots, p_K]^T$ on the $K$-dimensional simplex $\Delta^K$ is defined as:
\begin{equation}
    \text{Dir}(\bm{p} \mid \bm{\alpha}) = \frac{\Gamma(S)}{\prod_{c=1}^K \Gamma(\alpha_c)} \prod_{c=1}^K p_c^{\alpha_c - 1}
\end{equation}

\begin{remark}[Why not use ReLU?]
ReLU truncates negative logits to zero, resulting in zero gradients for those activations and permanently trapping evidence at $e_c = 0$, which freezes the Dirichlet state on the zero-evidence simplex. Therefore, the Softplus activation is strictly necessary to maintain a differentiable evidence manifold.
\end{remark}

\subsection{A.2 Proof of Vacuity Bounds}
By definition in Subjective Logic, Dirichlet vacuity is $u_e = \frac{K}{S}$, where $S = \sum_{c=1}^K \alpha_c$. Under our parameterization, $\alpha_c = \ln(1 + \exp(z_c)) + 1.0$. Since Softplus maps any real-valued logit to a positive value, $\alpha_c > 1.0$ for all $c$. Thus, $S = \sum_c \alpha_c > K$. Because $S > K$, $\frac{K}{S} < 1$. As $z_c \to -\infty$, $S \to K$, yielding supremum $u_e \to 1$. As $z_c \to \infty$, $S \to \infty$, yielding infimum $u_e \to 0$. Thus $u_e \in (0, 1]$.

\subsection{A.3 Expected Categorical Entropy Non-Negativity}
The expected categorical entropy is $u_a = \sum_{c=1}^K \frac{\alpha_c}{S} [\psi(S + 1) - \psi(\alpha_c + 1)]$. The digamma function $\psi(x)$ is strictly monotonically increasing for $x > 0$. Since $\alpha_c \ge 1$ and $S = \sum \alpha_i$, $S \ge \alpha_c$ for all $c$. Thus $\psi(S + 1) \ge \psi(\alpha_c + 1)$, meaning the term in brackets is non-negative. Thus $u_a \ge 0$.

\subsection{A.4 Asymmetric KL Divergence}
To shrink false evidence to the flat prior $1.0$, we adjust the Dirichlet parameters to $\tilde{\alpha}_c^{(n)} = y_c^{(n)} + (1 - y_c^{(n)}) \cdot \alpha_c^{(n)}$. The full KL divergence penalty $L_{\text{KL}}^{(n)}$ is computed as:
\begin{equation}
\begin{aligned}
    L_{\text{KL}}^{(n)} ={}& \ln \left( \frac{\Gamma\left(\sum_{c=1}^K \tilde{\alpha}_c^{(n)}\right)}{\Gamma(K) \prod_{c=1}^K \Gamma(\tilde{\alpha}_c^{(n)})} \right) \\
    &+ \sum_{c=1}^K (\tilde{\alpha}_c^{(n)} - 1)\left[\psi(\tilde{\alpha}_c^{(n)}) - \psi\!\left(\sum_{c=1}^K \tilde{\alpha}_c^{(n)}\right)\right]
\end{aligned}
\end{equation}

\section{Appendix B: Full GUDS-EDL Algorithm}
\begin{algorithm}[t]
\caption{Universal GUDS-EDL Training and Optimization Loop}
\label{alg:guds_edl}
\begin{algorithmic}[1]
\REQUIRE Dataset $\mathcal{D}$, model weights $\bm{W}^{(0)}$, latent scores $\bm{V}^{(0)} \leftarrow |\bm{W}^{(0)}|$, epochs $T$, warmup epochs $T_w$, structural update rate $\eta_{\text{struct}} = 0.02$.
\FOR{each epoch $t = 1 \dots T$}
    \IF{$t == T_w$}
        \STATE Freeze learned channel permutations into static index maps.
    \ENDIF
    \IF{$t \ge T_w$}
        \STATE Sample a proxy batch $(\bm{X}, \bm{y}) \sim \mathcal{D}$ for amortized topology computation.
        \STATE Compute structural gradients (Microglia and Astrocyte):
        \STATE \quad $\bm{C}_{\text{pruner}} \leftarrow \text{ReLU}\left( \bm{W} \odot \nabla_{\bm{W}} (u_a / (u_e + \epsilon)) \right)$
        \STATE \quad $\bm{G}_{\text{regrower}} \leftarrow \nabla_{\bm{W}} \left( \omega_y \cdot u_e \cdot \KL(\Dir(\bm{\alpha}) \parallel \Dir(\mathbf{1})) \right)$
        \STATE Update continuous latent scores $\bm{V}$ using $\bm{C}_{\text{pruner}}$ and $\bm{G}_{\text{regrower}}$.
        \STATE Apply layer-norm anti-crystallization noise if max growth is below threshold.
        \STATE Update binary mask $\bm{M}^{(l)}$ via Reshaped NVIDIA 2:4 top-2-of-4 selection on $\bm{V}^{(l)}$.
    \ELSE
        \STATE Set masks $\bm{M}^{(l)} \leftarrow \mathbf{1}$ (Dense Warmup).
    \ENDIF
    \FOR{each mini-batch step $s$ in epoch $t$}
        \STATE Sample batch $(\bm{X}, \bm{y}) \sim \mathcal{D}$.
        \STATE Compute learning rate $\eta_s$ and loss scale $S_{\text{loss}}(s)$ via warmup schedule.
        \STATE Forward pass: compute logits $\bm{z}$ using static masked weights $\bm{W}_{\text{eff}} = \bm{W} \odot \bm{M}$.
        \STATE Compute Dirichlet parameters: $\alpha_c \leftarrow \ln(1 + e^{z_c}) + 1.0$.
        \STATE Compute expected probabilities $\hat{p}_c = \alpha_c/S$ and uncertainties $u_e, u_a$.
        \STATE Compute Evidential Focal Loss (EFL) $L_{\text{EFL}}$.
        \STATE Backward pass: compute gradients $\nabla_{\bm{W}} L_{\text{scaled}}$.
        \STATE Update active weights: $\bm{W} \leftarrow \bm{W} - \eta_s \text{AdamW}(\nabla_{\bm{W}} L_{\text{scaled}} \odot \bm{M})$.
    \ENDFOR
\ENDFOR
\end{algorithmic}
\end{algorithm}


\section{Appendix C: Planned Ablation Protocol}
We plan a complete ablation suite to isolate the effects of the signed pruner, class-conditioned regrower, topology cache, anti-crystallization noise, asymmetric KL regularization, and bias-corrected calibration.


\section{Appendix D: Planned Generalization Benchmarks}
\subsection{D.1 CIFAR-100-LT Protocol}

\subsection{D.2 MVTec AD Image-Level Protocol}
\subsection{D.3 Hardware Profiling Protocol}
\subsection{D.4 Quality-Gated Failure Detection Protocol}

\section{Appendix E: Reproducibility Details}


\end{document}